\section{难度子项与权重依据的补充统计}

表~\ref{tab:appendix_weight_evidence} 汇总本文已使用结果文件中的代表性统计，数据来自 MBPP 上的离线路由证据与特征实验报告：一类是子项与结构分档/运行标签的 Spearman 相关，另一类是学习式路由基线在不同特征块上的交叉验证 AUC。该表用于支撑第~\ref{sec:mechverify} 节关于结构分主要承担预算排序先验、而非运行失败分类器的判断。

\begin{table}[H]
\centering
\scriptsize
\caption{难度子项与权重依据的补充统计}
\label{tab:appendix_weight_evidence}
\begin{tabularx}{\linewidth}{>{\raggedright\arraybackslash}X c c c}
\toprule
统计项 & 指标 & 数值 & $n$ \\
\midrule
\makecell[l]{$\rho(\texttt{structural\_hard},$\\$\texttt{test\_score})$} & Spearman & 0.5696 & 427 \\
\makecell[l]{$\rho(\texttt{direct\_fail},$\\$\texttt{test\_score})$} & Spearman & -0.0337 & 383 \\
\makecell[l]{$\rho(\texttt{direct\_fail},$\\$\texttt{boundary\_score})$} & Spearman & 0.0483 & 383 \\
\makecell[l]{$\rho(\texttt{direct\_fail},$\\$\texttt{extent\_score})$} & Spearman & 0.0669 & 383 \\
\makecell[l]{\texttt{need\_reroute}\\（v2 五子分）} & LR 5-fold AUC & 0.450 $\pm$ 0.103 & 383 \\
\makecell[l]{\texttt{need\_reroute}\\（扩展静态）} & LR 5-fold AUC & 0.520 $\pm$ 0.108 & 383 \\
\bottomrule
\end{tabularx}
\end{table}

该补充表支持本文的方法定位：结构难度与测试复杂度的相关系数为 0.5696，说明测试复杂度能够进入结构先验；而首轮失败与各静态子分的相关性接近 0，LR 5 折 AUC 也仅从 0.450 提升到 0.520。由此可见，难度子项更适合作为\textbf{可解释的预算分配先验}，而非直接充当运行失败分类器。主结论仍以端到端通过率、调用与令牌前沿为准。
